\section*{अध्याय \ref{ch:b1:plate-kinematics} --- गोले पर प्लेट गतिकी}

\begin{solution}{exo:b1:plate-kinematics:1}
(क) $\qty{0.210}{\degree\per\Myr} = 0.210\times\pi/180 = \qty{3.665e-3}{rad}$
प्रति दस लाख वर्ष, अर्थात् \qty{3.665e-9}{rad.yr^{-1}}; $\qty{1}{yr} = \qty{3.156e7}{s}$ के साथ, \qty{1.161e-16}{rad.s^{-1}}। (ख) $360/0.210 =
\qty{1714}{\Myr}$, पृथ्वी की आयु के एक-तिहाई से अधिक। (ग) $\omega R =
\num{3.665e-9}\times\qty{6.371e9}{mm} = \qty{23.35}{mm.yr^{-1}}$।
\end{solution}

\begin{solution}{exo:b1:plate-kinematics:2}
$v = \omega R\sin\Delta = 23.35\sin\Delta$: \ang{30} पर \qty{11.7}{mm.yr^{-1}},
\ang{60} पर \qty{20.2}{mm.yr^{-1}} और \ang{90} पर \qty{23.35}{mm.yr^{-1}}। प्रसार वहाँ सबसे तेज़ है जहाँ सीमा ध्रुव से सबसे दूर है, \ang{90} तक: यूरेशिया और उत्तरी अमेरिका के बीच, सीमा के दक्षिणी सिरे पर, अज़ोरेस के निकट, ध्रुव से \ang{78} दूर।
\end{solution}

\begin{solution}{exo:b1:plate-kinematics:3}
उपकालखंड $1.072 - 0.988 = \qty{0.084}{\Myr}$ तक चला, और \unit{mm.yr^{-1}} में अर्ध दर \unit{km\per\Myr} में भी एक दर है। (क) \qty{50}{mm.yr^{-1}} पर $50\times0.084 =
\qty{4.2}{km}$; \qty{10}{mm.yr^{-1}} पर $10\times0.084 = \qty{0.84}{km}$। (ख) धारी अक्ष के आर-पार फैलती है:
$2\times10\times0.781 = \qty{15.6}{km}$।
\end{solution}

\begin{solution}{exo:b1:plate-kinematics:4}
अर्ध दरें $66/2.581 = \qty{25.6}{mm.yr^{-1}}$ पश्चिम की ओर और
$64/2.581 = \qty{24.8}{mm.yr^{-1}}$ पूर्व की ओर हैं; पूर्ण दर
$130/2.581 = \qty{50.4}{mm.yr^{-1}}$ है। प्रसार लगभग सममित है: पश्चिमी पार्श्व लगभग \qty{3}{\percent} तेज़ बढ़ा है।
\end{solution}

\begin{solution}{exo:b1:plate-kinematics:5}
$\cos\Delta = \sin\ang{61.8}\sin\ang{85} + \cos\ang{61.8}\cos\ang{85}
\cos(\ang{0} - \ang{139.6}) = 0.8813\times0.9962 + 0.4726\times0.08716\times
(-0.7615) = 0.87795 - 0.03137 = 0.84658$, अतः $\Delta = \ang{32.16}$ और
$v = 23.35\sin\ang{32.16} = \qty{12.4}{mm.yr^{-1}}$। यह मापी गई \qtyrange{11}{12}{mm.yr^{-1}} से थोड़ा ऊपर है, पर \omterm{def:b1:plate-kinematics:euler-pole}{यूलर सदिश} 453 दरों पर एक एकल सर्वोत्तम फ़िट है, जिनसे दरें औसतन \qty{0.7}{mm.yr^{-1}} विचलित होती हैं, और धीमे गाक्कल कटक के निकट की व्यक्तिगत दरें सबसे कम सटीक में से हैं।
\end{solution}

\begin{solution}{exo:b1:plate-kinematics:6}
(क) $\cos\Delta = \sin\ang{61.8}\sin\ang{39.44} + \cos\ang{61.8}\cos\ang{39.44}
\cos(\ang{169.43}) = 0.20111$: $\Delta = \ang{78.40}$। (ख) $v = 23.35
\sin\ang{78.40} = \qty{22.87}{mm.yr^{-1}}$, जो संधि के ठीक उत्तर में मापी गई \qty{23}{mm.yr^{-1}} है। (ग) विधि के सूत्र से, अंश
$\sin\ang{169.43}\cos\ang{61.8} = 0.08668$ है और हर
$\cos\ang{39.44}\sin\ang{61.8} - \sin\ang{39.44}\cos\ang{61.8}\cos\ang{169.43}
= 0.97573$: ध्रुव $\beta = \ang{5.08}$ पर दिखता है, लगभग ठीक उत्तर, और यूरेशिया $\ang{95.1}$ की ओर चलता है। (घ) रूपांतर सापेक्ष गति के समांतर हैं: वे लगभग \ang{095} की प्रवृत्ति रखते हैं, लगभग पूर्व-पश्चिम।
\end{solution}

\begin{solution}{exo:b1:plate-kinematics:7}
(क) $\vec\omega = \omega(\cos\varphi_p\cos\lambda_p, \cos\varphi_p\sin\lambda_p,
\sin\varphi_p)$: $\vec\omega_{\mathrm{EU\,NA}} = (-0.0756, 0.0643, 0.1851)$ और
$\vec\omega_{\mathrm{NB\,NA}} = (0.0333, 0.0282, 0.2289)$, \unit{\degree\per\Myr} में। (ख) $\vec\omega_{\mathrm{NB\,EU}} =
\vec\omega_{\mathrm{NB\,NA}} - \vec\omega_{\mathrm{EU\,NA}} = (0.1089, -0.0361,
0.0438)$, जिसकी लंबाई
\[
  \sqrt{0.011859 + 0.001303 + 0.001918} = \qty{0.1228}{\degree\per\Myr}.
\]
उसका ध्रुव \omterm{def:b1:spherical-astronomy:geographic}{अक्षांश}
$\arcsin(0.0438/0.1228) = \ang{20.9}$\,N और \omterm{def:b1:spherical-astronomy:geographic}{देशांतर}
$\arctan(-0.0361/0.1089) = \ang{-18.3}$, अर्थात् \ang{18.3}\,W है ($x$ घटक धनात्मक होने से कोई \ang{180} नहीं जोड़ा जाता)। (ग) नूबिया और यूरेशिया के बीच की सीमा के प्रत्यक्ष माप कम हैं: अज़ोरेस और जिब्राल्टर के बीच, केवल ग्लोरिया भ्रंश के दिगंश। अटलांटिक के आर-पार दोनों प्रसार-गतियाँ सैकड़ों दरों से मापी गई हैं, और उनका अंतर उनकी सटीकता विरासत में पाता है।
\end{solution}

\begin{solution}{exo:b1:plate-kinematics:8}
(क) $\vec e_z$ के परितः \ang{90}: $\cos\ang{90} = 0$, $\vec e_z\cdot\vec e_x = 0$
और $\vec e_z\times\vec e_x = \vec e_y$, अतः $\vec e_x$ $\vec e_y$ पर जाता है। फिर $\vec e_x$ के परितः \ang{90}: $\vec e_x\cdot\vec e_y = 0$ और
$\vec e_x\times\vec e_y = \vec e_z$, अतः $\vec e_y$ $\vec e_z$ पर जाता है: बिंदु उत्तरी ध्रुव पर समाप्त होता है। (ख) पहले $\vec e_x$ के परितः: $\vec e_x$ अक्ष पर है,
$\vec e_x\times\vec e_x = \vec 0$ और $(1 - \cos\ang{90})(\vec e_x\cdot\vec e_x)
\vec e_x = \vec e_x$, अतः वह नहीं चलता; फिर $\vec e_z$ के परितः वह $\vec e_y$ पर जाता है, \ang{0}\,N, \ang{90}\,E पर। (ग) परिणाम भिन्न हैं: \omterm{def:b1:plate-kinematics:finite-rotation}{परिमित घूर्णन} क्रम-विनिमेय नहीं होते। क्रमिक पुनर्रचनाएँ आव्यूहों को सही क्रम में गुणा करके संयोजित होनी चाहिए, कभी कोणों या ध्रुवों को जोड़कर नहीं।
\end{solution}

\begin{solution}{exo:b1:plate-kinematics:9}
(क) वेग-अंतरिक्ष में $a$ को मूलबिंदु पर रखिए, $b$ को \qty{30}{mm.yr^{-1}} पूर्व की ओर और $c$ को \qty{40}{mm.yr^{-1}} दक्षिण की ओर। (ख) $\vec v_{CB} =
\vec v_{CA} - \vec v_{BA} = (0, -40) - (30, 0) = (-30, -40)$ (पूर्व, उत्तर घटक):
\qty{50}{mm.yr^{-1}} दिगंश की ओर
$\arctan\bigl((-30)/(-40)\bigr) + \ang{180} = \ang{216.9}$, दक्षिण-पश्चिम। (ग) कटक इस दिशा के लंबवत है: उसकी प्रवृत्ति $\ang{216.9} - \ang{90} = \ang{126.9}$ है, अर्थात् उत्तर-पश्चिम–दक्षिण-पूर्व।
\end{solution}

\begin{solution}{exo:b1:plate-kinematics:10}
(क) रॉड्रिग्ज़ के आव्यूह में, तत्समक $3\cos\alpha$ का योगदान देता है, आव्यूह
$\vec n\,\vec n^{\mathsf T}$ का ट्रेस $n_x^2 + n_y^2 + n_z^2 = 1$ है, जो $1 - \cos\alpha$ से गुणा होता है, और प्रतिसममित आव्यूह कुछ नहीं: $\operatorname{tr}R =
3\cos\alpha + 1 - \cos\alpha = 1 + 2\cos\alpha$। (ख) यहाँ $\operatorname{tr}R = 0$,
अतः $\cos\alpha = -1/2$ और $\alpha = \ang{120}$। अक्ष $R\vec n =
\vec n$ को संतुष्ट करता है: $R(x, y, z) = (z, x, y)$, अतः $x = y = z$ और $\vec n =
(1, 1, 1)/\sqrt3$, \omterm{def:b1:spherical-astronomy:geographic}{अक्षांश} $\arcsin(1/\sqrt3) = \ang{35.26}$\,N और
\omterm{def:b1:spherical-astronomy:geographic}{देशांतर} \ang{45}\,E पर। (ग) $R\vec e_x = \vec e_y$, $R\vec e_y = \vec e_z$,
$R\vec e_z = \vec e_x$: तीनों अक्ष क्रमचयित होते हैं, जैसा विकर्ण के परितः एक-तिहाई चक्कर करता है। रॉड्रिग्ज़ का सूत्र $\alpha = +\ang{120}$ के साथ, $\vec n$ के परितः, सचमुच
$\vec e_x\mapsto -\tfrac12\vec e_x + \tfrac12(1, 1, 1) +
\tfrac{\sqrt3}{2}\cdot\tfrac{1}{\sqrt3}(0, 1, -1) = (0, 1, 0)$ देता है, अतः घूर्णन \ang{35.26}\,N, \ang{45}\,E के ऊपर से देखने पर वामावर्त है।
\end{solution}

\begin{solution}{exo:b1:plate-kinematics:11}
(क) एक रूपांतर पर सापेक्ष वेग $\vec\omega\times\vec r$ भ्रंश के अनुदिश है, अतः $\vec t$ के समांतर; वह $\vec\omega$ के लंबवत है, इसलिए
$\vec t\cdot\vec\omega = 0$। केंद्र से होकर $\vec t$ के लंबवत तल में $\vec r$ (क्योंकि $\vec t$ क्षैतिज है) और $\vec\omega$ हैं: वह बिंदु से होकर जाने वाले \omterm{def:b1:spherical-astronomy:spherical-triangle}{महावृत्त} का तल है जो भ्रंश के लंबवत है, और ध्रुव उस वृत्त पर है। ऐसे दो तल $\vec t_1\times\vec t_2$ की रेखा के अनुदिश मिलते हैं। (ख) $\vec t_1 = (0.5436, 0.8367, -0.0673)$ संधि पर और $\vec t_2 = (0.4942, 0.8639, -0.0968)$ आइसलैंड में, अतः
$\vec t_1\times\vec t_2 = (-0.02285, 0.01936, 0.05612)$, लंबाई $0.06362$ की। इसका \omterm{def:b1:spherical-astronomy:geographic}{अक्षांश} $\arcsin(0.05612/0.06362) = \ang{61.9}$\,N और इसका \omterm{def:b1:spherical-astronomy:geographic}{देशांतर}
$\ang{180} - \arctan(0.01936/0.02285) = \ang{139.7}$\,E ($x$ घटक ऋणात्मक होने से): अध्याय का ध्रुव। (ग) बुरी तरह: दोनों \omterm{def:b1:spherical-astronomy:spherical-triangle}{महावृत्त} एक छोटे कोण पर कटते हैं, क्योंकि सीमा के \ang{26} पर प्रवृत्तियाँ केवल \ang{8} भिन्न हैं। किसी एक प्रवृत्ति को एक डिग्री बदलने से ध्रुव लगभग \ang{15} चलता है। वास्तविक निर्धारण कई रूपांतरों को प्रसार-दरों के साथ जोड़ते हैं, जिनकी सीमा के साथ भिन्नता ध्रुव से कोणीय दूरी तय करती है।
\end{solution}

\begin{solution}{exo:b1:plate-kinematics:12}
(क) $(94 - 4)/3 = \qty{30}{mm.yr^{-1}}$। (ख) एक वर्ष के अंतर पर अंतर हैं
$24 - 4 = 20$, $25 - 5 = 20$, $26 - 6 = 20$, $65 - 15 = 50$, $74 - 24 = 50$,
$75 - 25 = 50$, $76 - 26 = 50$, $85 - 65 = 20$ और $94 - 74 = 20$ (mm)। छाँटने पर:
पाँच मान 20 के और चार 50 के; माध्यिका \qty{20}{mm.yr^{-1}} है। (ग) एक वर्ष के अंतर पर दो स्थितियाँ मौसमी दोलन के एक ही चरण पर होती हैं, जो उनके अंतर में निरस्त हो जाता है। एक छलांग केवल उन जोड़ों को प्रभावित करती है जो उसे पार करते हैं, यहाँ नौ में से चार; जब तक आधे से कम जोड़े किसी छलांग को पार करते हैं, माध्यिका उसे अनदेखा करती है, जबकि (क) का अंत-बिंदु आकलन पूरी छलांग को गति मानता है और \qty{50}{\percent} अधिक निकलता है। विधि तब विफल होती है जब छलांगें इतनी बार आएँ कि अधिकांश एक-वर्षीय जोड़ों में एक हो, या जब स्टेशन अस्थिर रूप से चले, उदाहरणार्थ एक बड़े भूकंप के बाद।
\end{solution}

\begin{solution}{pb:b1:plate-kinematics:1}
\textbf{1.} $\omega = 0.210\times(\pi/180)/\num{e6} = \qty{3.665e-9}{rad.yr^{-1}}$;
$\omega R = \num{3.665e-9}\times\qty{6.371e9}{mm} = \qty{23.35}{mm.yr^{-1}}$।

\textbf{2.} $\cos\Delta = \sin\ang{61.8}\sin\ang{64.5} + \cos\ang{61.8}
\cos\ang{64.5}\cos(\ang{-157.6}) = 0.79546 - 0.18809 = 0.60737$, अतः
$\Delta = \ang{52.60}$।

\textbf{3.} $v = 23.35\sin\ang{52.60} = 23.35\times0.7944 =
\qty{18.55}{mm.yr^{-1}}$।

\textbf{4.} अंश $\sin\ang{157.6}\cos\ang{61.8} = 0.18011$ है और हर
\[
  \cos\ang{64.5}\sin\ang{61.8} - \sin\ang{64.5}\cos\ang{61.8}\cos\ang{157.6} =
  0.37940 + 0.39436 = 0.77376;
\]
दोनों धनात्मक हैं, अतः $\beta = \arctan 0.2328 =
\ang{13.1}$। यूरेशिया $\ang{13.1} + \ang{90} = \ang{103.1}$ की ओर चलता है।

\textbf{5.} पूर्व: $18.55\sin\ang{103.1} = \qty{18.07}{mm.yr^{-1}}$; उत्तर:
$18.55\cos\ang{103.1} = \qty{-4.20}{mm.yr^{-1}}$, अर्थात् दक्षिण की ओर \qty{4.20}{mm.yr^{-1}}।

\textbf{6.} $\cos\Delta = 0.8851$, $\Delta = \ang{27.74}$, $v = 23.35
\sin\ang{27.74} = \qty{10.9}{mm.yr^{-1}}$। दोनों प्लेटें पूर्वी साइबेरिया के एक ध्रुव के परितः घूमती हैं: गाक्कल कटक उसके निकट है, और चाल $\sin\Delta$ के साथ घटती है, यद्यपि प्लेटें वही हैं।

\textbf{7.} पश्चिम: $94/9.779 = \qty{9.61}{mm.yr^{-1}}$; पूर्व: $97/9.779 =
\qty{9.92}{mm.yr^{-1}}$।

\textbf{8.} $(94 + 97)/9.779 = 191/9.779 = \qty{19.53}{mm.yr^{-1}}$।

\textbf{9.} C2An से: $54.1/2.581 = \qty{20.96}{mm.yr^{-1}}$; C1n के आधार से:
$19.8/0.781 = \qty{25.35}{mm.yr^{-1}}$। वे सहमत नहीं हैं: सीमा जितनी युवा, दर उतनी ऊँची। दर का एक वास्तविक परिवर्तन संभव है, पर दूरियों में एक स्थिर अतिरिक्त, जो सभी आयुओं पर समान हो, ठीक यही करेगा।

\textbf{10.} पश्चिम और पूर्व दूरियों का अनुपात 0.96 से 0.98 है: पूर्वी पार्श्व कुछ प्रतिशत अधिक बढ़ा है। पूर्ण दर केवल दो दूरियों के योग पर निर्भर करती है, अतः असममिति उसे नहीं बदलती।

\textbf{11.} $(191 - 5.0)/9.779 = \qty{19.02}{mm.yr^{-1}}$;
$(54.1 - 5.0)/2.581 = \qty{19.02}{mm.yr^{-1}}$; $(19.8 - 5.0)/0.781 =
\qty{18.95}{mm.yr^{-1}}$। तीनों अब \qty{0.1}{mm.yr^{-1}} के भीतर सहमत हैं। वही \qty{5}{km} सबसे युवे जोड़े के \qty{19.8}{km} का एक-चौथाई है, परंतु सबसे पुराने के \qty{191}{km} का \qty{3}{\percent} से कम; पुरानी सीमाएँ अधिक विश्वसनीय दरें देती हैं।

\textbf{12.} सुधारी अर्ध दर $186/9.779/2 = \qty{9.51}{mm.yr^{-1}}$ है,
और $9.51\times0.084 = \qty{0.80}{km}$, वह \qty{2.5}{km} का एक-तिहाई जिसके द्वारा हर सीमा विस्थापित है। इतनी संकरी धारी का धुंधला हो जाना अपेक्षित होता। कि वह रेक्यनीस कटक पर पहचानी जाती है, जबकि गाक्कल कटक पर, जिसमें मैग्मा की कमी है, अनुपस्थित है, सुझाता है कि मैग्मा की प्रचुर आपूर्ति संकरी धारियों में भी क्षेत्र को विश्वसनीय रूप से दर्ज करती है।

\textbf{13.} सुधारी दरें, \qty{18.95} से \qty{19.02}{mm.yr^{-1}}, प्लेट सदिश की \qty{19.0}{mm.yr^{-1}} से सहमत हैं; सुधार के बिना वे उससे 3 से 33 प्रतिशत अधिक थीं।

\textbf{14.} $\cos\Delta = \sin\ang{64.139}\sin\ang{59.018} +
\cos\ang{64.139}\cos\ang{59.018}\cos\ang{27.554} = 0.97054$, अतः $\Delta =
\ang{13.94}$; जीवा $2R\sin(\Delta/2) = 12742\times\sin\ang{6.97} =
\qty{1546}{km}$ है।

\textbf{15.} $1000/17.5 = 57$ वर्ष।

\textbf{16.} दो दिशाएँ $\ang{123.1} - \ang{114.8} =
\ang{8.3}$ से भिन्न हैं: \omterm{def:b1:spherical-astronomy:spherical-triangle}{महावृत्त} के अनुदिश क्षैतिज घटक, यानी
$18.9\cos\ang{8.3} = \qty{18.70}{mm.yr^{-1}}$ है, और जीवा के अनुदिश
$18.70\cos\ang{6.97} = \qty{18.56}{mm.yr^{-1}}$, लगभग \qty{18.6}{mm.yr^{-1}}।

\textbf{17.} जीवा \qty{17.5}{mm.yr^{-1}} से लंबी होती है, जो भविष्यवाणित \qty{18.6}{mm.yr^{-1}} से \qty{6}{\percent} कम है। दोनों स्टेशन प्लेट सीमा से दूर हैं, इसलिए अंतर स्थानीय नहीं है। या तो प्लेटें आज पिछले तीन मिलियन वर्षों के औसत से थोड़ी धीमी चल रही हैं, या दो आकलनों में से एक में एक छोटी व्यवस्थित त्रुटि है; अकेले माप यह नहीं बता सकते। उपग्रह स्थिति-निर्धारण अकेले से परिकलित \omterm{def:b1:plate-kinematics:euler-pole}{यूलर सदिश} इस जीवा के लिए \qty{17.1}{mm.yr^{-1}} की भविष्यवाणी करता है, जो मापे गए मान के निकट है।

\textbf{18.} पूर्व: $21.12 - (-1.07) = \qty{22.19}{mm.yr^{-1}}$; उत्तर:
$-6.97 - (-0.97) = \qty{-6.00}{mm.yr^{-1}}$। चाल $\sqrt{22.19^2 + 6.00^2} =
\qty{23.0}{mm.yr^{-1}}$, की ओर $\ang{180} - \arctan(22.19/6.00) =
\ang{105.1}$।

\textbf{19.} दिशा प्लेट सदिश के \ang{103} के निकट है, पर चाल उसकी \qty{18.6}{mm.yr^{-1}} को एक-चौथाई से अधिक करती है। दोनों स्टेशन उस क्षेत्र के भीतर हैं जहाँ सीमा विकृत होती है, और होफ़्न, सबसे बड़ी हिम-टोपी के निकट, बर्फ़ के पतले होने पर \omterm{def:b1:layered-earth:layers}{भूपर्पटी} के उछाल के साथ चलता है: ऊपर की ओर \qty{12.9}{mm.yr^{-1}}, और बर्फ़ से दूर, बगल की ओर। सीमा-क्षेत्र के भीतर के स्टेशन प्लेट गति और स्थानीय विकृति दोनों मापते हैं।

\textbf{20.} $0.210\times10 = \ang{2.10}$।

\textbf{21.} बिंदु त्रिज्या $R\sin\Delta$ के एक लघुवृत्त पर चलता है:
$6371\times0.7944\times(2.10\times\pi/180) = \qty{185.5}{km}$। C5n के युवा किनारे के बाद बनी \omterm{def:b1:layered-earth:layers}{भूपर्पटी} की सुधारी चौड़ाई, \qty{9.8}{Myr} में \qty{186}{km}, लगभग वही है: गति दस लाख वर्षों तक स्थिर रही है।

\textbf{22.} घटक (पूर्व, उत्तर): यूरेशिया के लिए
$(22.9\sin\ang{95.1}, 22.9\cos\ang{95.1}) =$ $(22.81, -2.04)$, नूबिया के लिए:
\[ (19.2\sin\ang{103.8}, 19.2\cos\ang{103.8}) = (18.65, -4.58) \]
नूबिया के सापेक्ष यूरेशिया: $(4.16, 2.54)$, अर्थात् \qty{4.9}{mm.yr^{-1}} की ओर
$\arctan(4.16/2.54) = \ang{58.6}$, पूर्व-उत्तर-पूर्व।

\textbf{23.} तीन कटकों की एक संधि वेगों की परवाह किए बिना \omterm{def:b1:plate-kinematics:triple-junction}{स्थायी} है। प्रश्न 22 का वेग यूरेशिया और नूबिया को पूर्व-उत्तर-पूर्व से पश्चिम-दक्षिण-पश्चिम की ओर दूर करता है: संधि पर उनकी सीमा धीरे खुलनी चाहिए; यह टर्सीरा दरार है, जिसकी पूर्व-दक्षिण-पूर्व प्रवृत्ति खुलाव को तिरछा बनाती है।

\textbf{24.} भाग I: \qty{18.6}{mm.yr^{-1}} आइसलैंड के केंद्र में
(\qty{19.0}{mm.yr^{-1}} भाग II की परिच्छेदिका पर)। भाग II: सुधार के बाद 18.95 से
\qty{19.02}{mm.yr^{-1}}। भाग III: रेक्याविक–स्टैवेंगर जीवा \qty{17.5}{mm.yr^{-1}} से लंबी होती है, जहाँ प्लेट सदिश \qty{18.6}{mm.yr^{-1}} की भविष्यवाणी करता है; आइसलैंड के भीतर की जोड़ी \qty{23.0}{mm.yr^{-1}} देती है, जो प्लेट गति नहीं है। प्लेट सदिश और धारियाँ सहमत हैं; उपग्रह \qty{6}{\percent} भिन्न हैं। धारियाँ दस मिलियन वर्षों तक औसत करती हैं, प्लेट सदिश \qty{3.16}{Myr} पर; उपग्रह पिछले बीस से तीस वर्षों की गति मापते हैं।

\textbf{25.} प्लेट सदिश आइसलैंड के आर-पार \qty{18.55}{mm.yr^{-1}} देता है, और उसके बगल में सुधारी धारियाँ \qty{19.0}{mm.yr^{-1}}: एक मिलीमीटर प्रति वर्ष तक पूर्णांकित, आइसलैंड के आर-पार \omterm{def:b1:plate-kinematics:spreading-rate}{पूर्ण प्रसार-दर}
\textbf{\qty[text-series-to-math=true]{19}{mm.yr^{-1}}} है।
\end{solution}
